\documentclass[11pt]{article}
\usepackage[margin=1in]{geometry}
\usepackage{amsmath,amssymb,amsthm}
\usepackage{xurl}
\usepackage{hyperref}
\sloppy
\let\oldus\_ \renewcommand{\_}{\oldus\allowbreak}
\newtheorem{theorem}{Theorem}
\newtheorem{lemma}[theorem]{Lemma}
\theoremstyle{definition}
\newtheorem{definition}[theorem]{Definition}

\title{Disproof of conjecture 00000009682:\\ two algebraically independent U-numbers exist}
\author{Jackmeson1}
\date{2026-10-05}

\begin{document}
\maketitle

\section{The conjecture and the clause refuted}

\textbf{Statement (English, verbatim).} Definition: Mahler's classification: transcendental
numbers by approximation type into classes A, S, T, U; the measure and cardinality structure of
the classes. Conjecture: Lebesgue-almost every transcendental number is an S number, and the
U-number set has measure zero but continuum cardinality; the dimension of algebraically
independent subsets of U numbers is exactly 1, with no two-dimensional U families. (Mahler class
measure structure)

The Chinese text states the same; its last clause says literally that the algebraically
independent subsets of the U-numbers have dimension exactly 1 and that there is no
two-dimensional U family.

\textbf{Reading.} The conjecture is a conjunction of
(C1) almost every transcendental number is an S-number;
(C2) the set of U-numbers is Lebesgue-null and has cardinality $\mathfrak c$; and
(C3) ``the dimension of algebraically independent subsets of U numbers is exactly 1, with no
two-dimensional U families''. We read (C3) as the conjunction of
\begin{itemize}
\item[(C3a)] the supremum of $|s|$ over the subsets $s$ of the U-numbers that are
  algebraically independent over $\mathbb Q$ equals $1$ (dimension = transcendence degree), and
\item[(C3b)] there is no two-element subset of the U-numbers that is algebraically independent
  over $\mathbb Q$.
\end{itemize}
We refute (C3a) and (C3b) separately, for real numbers and for complex numbers
(Lebesgue measure on $\mathbb R$, resp.\ $\mathbb C$). This refutes the conjunction in both
settings. We do not use (C1) or (C2).

\textbf{Readings not covered.} (i) A reading of ``dimension'' as Hausdorff dimension of a set
is not treated; the text attaches ``dimension'' to algebraically independent subsets, so we use
the algebraic (transcendence-degree) reading. (ii) Our two witnesses are U-numbers of degree $1$
(both have $\omega(\cdot,1)=\infty$). A reading restricted to U-numbers of one fixed degree
$n\ge 2$ is not refuted here. (iii) The class ``A'' in the Definition line plays no role; we
only use the U- and S-classes.

\section{Definitions (Wikipedia, ``Transcendental number theory'', section ``Mahler's
classification'')}

Source: \url{https://en.wikipedia.org/wiki/Transcendental_number_theory}. Quoted:
``Next consider the values of polynomials at a complex number x, when these polynomials have
integer coefficients, degree at most n, and height at most H, with n, H being positive
integers.'' $m(x,n,H)$ is ``the minimum non-zero absolute value such polynomials take at'' $x$;
$\omega(x,n,H)=-\log m(x,n,H)/(n\log H)$ and $\omega(x,n)=\limsup_{H\to\infty}\omega(x,n,H)$.
``Suppose this is infinite for some minimum positive integer n. A complex number x in this case
is called a U number of degree n.'' ``If the $\omega$(x, n) are bounded, then $\omega$(x) is
finite, and x is called an S number.'' The same article: ``Kurt Mahler in 1932 partitioned the
transcendental numbers into 3 classes, called S, T, and U.''

\begin{definition}[Lean: \texttt{values}, \texttt{mahlerM}, \texttt{mahlerOmegaH},
\texttt{mahlerOmega}]
For $x\in\mathbb C$ and $n,H\in\mathbb N$, let $V(x,n,H)$ be the set of non-zero values
$|P(x)|$ with $P\in\mathbb Z[X]$, $\deg P\le n$ and height $\le H$, i.e.\ $|a_i|\le H$ for
every coefficient $a_i$ of $P$. Put $m(x,n,H)=\inf V(x,n,H)$,
$\omega(x,n,H)=-\log m(x,n,H)/(n\log H)$, and
$\omega(x,n)=\limsup_{H\to\infty}\omega(x,n,H)$ in the extended reals
$[-\infty,\infty]$, where $H$ runs through $\mathbb N$.
\end{definition}

\begin{definition}[Lean: \texttt{IsUNumber}, \texttt{IsSNumber}]
$x\in\mathbb C$ is a \emph{U-number} if $x$ is transcendental over $\mathbb Q$ and
$\omega(x,n)=+\infty$ for some integer $n\ge1$. It is an \emph{S-number} if it is
transcendental and there is a real $C$ with $\omega(x,n)\le C$ for all $n\ge1$. A real number
is a U-number (S-number) when it is one as a complex number; the polynomial values are the same.
\end{definition}

\textbf{Conventions.} $\deg P\le n$ is \texttt{P.natDegree <= n}. Mathlib's $\mathbb N$
contains $0$, so $n\ge1$ is required explicitly. In Lean, $\log 0=0$, $\log 1=0$ and $t/0=0$;
these conventions affect $\omega(x,n,H)$ only for $H\le1$, which the $\limsup$ ignores. Since
$V(x,n,H)$ is finite (Lemma~\ref{lem:fin}), $\inf$ is the minimum whenever $V\ne\emptyset$.
\emph{Algebraic independence} over $\mathbb Q$ is Mathlib's \texttt{AlgebraicIndependent}: the
evaluation map from the multivariate polynomial ring is injective. The same Wikipedia article:
``A set of numbers \{$\alpha_1$, $\alpha_2$, \ldots, $\alpha_n$\} is called algebraically
independent over a field K if there is no non-zero polynomial P in n variables with coefficients
in K such that P($\alpha_1$, $\alpha_2$, \ldots, $\alpha_n$) = 0.'' A \emph{Liouville number}
is Mathlib's \texttt{Liouville x}: for every $n$ there are integers $a$, $b>1$ with
$0<|x-a/b|<b^{-n}$.

\textbf{The formal conjecture} (Lean: \texttt{ConjectureComplex}; \texttt{ConjectureReal} is the
same over $\mathbb R$):
(C1) $\forall^{\text{a.e.}}x$, transcendental $\Rightarrow$ S-number;
(C2) the U-set has volume $0$ and $\#\{x : \text{U-number}\}=\mathfrak c$;
(C3a) $\sup\{\,|s| : s\subseteq\mathrm U,\ s \text{ alg.\ indep.\ over }\mathbb Q\,\}=1$
(\texttt{Set.encard}, in $\mathbb N\cup\{\infty\}$); and
(C3b) no such $s$ has $|s|=2$.

\section{Proof}

\begin{lemma}[\texttt{finite\_polys}, \texttt{values\_finite}, \texttt{mahlerM\_spec}]
\label{lem:fin}
The integer polynomials of degree $\le n$ and height $\le H$ form a finite set. Hence if such a
$P$ has $P(x)\ne0$, then $0<m(x,n,H)\le|P(x)|$.
\end{lemma}
\begin{proof}
$P\mapsto(a_0,\dots,a_n)$ is injective on this set, with values in the finite box
$[-H,H]^{n+1}$. So $V(x,n,H)$ is finite and non-empty, and its infimum is one of its elements,
which are positive. The infimum is at most $|P(x)|\in V(x,n,H)$.
\end{proof}

\begin{theorem}[\texttt{liouville\_mahlerOmega\_one}, \texttt{liouville\_isUNumber}]
If $x\in\mathbb R$ is a Liouville number, then $\omega(x,1)=+\infty$, and $x$ is a U-number.
\end{theorem}
\begin{proof}
Fix $w\in\mathbb R$ and $N$. Let $k=\lceil 2w\rceil_+ +1$, so $k-1\ge 2w$ and $k-1\ge0$.
By Mathlib's \texttt{Liouville.frequently\_exists\_num}, there are $b\ge\max(N,2,\lceil|x|\rceil+1)$
and $a\in\mathbb Z$ with $0<|x-a/b|<b^{-k}$. Let $P=bX-a$ and $H=\max(b,|a|)\ge N$. Then
$\deg P\le1$, $P$ has height $\le H$, and $0<|P(x)|=b|x-a/b|<b^{1-k}$. By
Lemma~\ref{lem:fin}, $0<m(x,1,H)\le|P(x)|$, so $-\log m(x,1,H)>(k-1)\log b$. Also
$|a|\le b|x|+|bx-a|\le b(|x|+1)$, so $H\le b(|x|+1)\le b^2$ (as $|x|+1\le b$). Hence
$0<\log H\le2\log b$, and
\[
\omega(x,1,H)=\frac{-\log m(x,1,H)}{\log H}\ \ge\ \frac{(k-1)\log b}{2\log b}\ \ge\ w
\]
(the case $w\le0$ is immediate). So $\omega(x,1,H)\ge w$ for arbitrarily large $H$, and
$\omega(x,1)\ge w$ for every real $w$; thus $\omega(x,1)=+\infty$. Liouville numbers are
transcendental (Mathlib \texttt{Liouville.transcendental}, over $\mathbb Z$, hence over
$\mathbb Q$), so $x$ is a U-number. This matches the source: ``Clearly the Liouville numbers are
a subset of the U numbers.''
\end{proof}

\begin{lemma}[\texttt{eventually\_transcendental\_adjoin}]
Let $A$ be $\mathbb R$ or $\mathbb C$ (in Lean: any field extension of $\mathbb R$) and
$x\in\mathbb R$. Then residually many $y\in\mathbb R$ (all $y$ outside a meagre set) are
transcendental over the $\mathbb Q$-algebra $K=\mathbb Q[x]\subseteq A$.
\end{lemma}
\begin{proof}
$K$ is countable, since $\#\mathbb Q[x]\le\#\mathbb Q+\#\{x\}+\aleph_0$
(\texttt{Algebra.cardinalMk\_adjoin\_le}). The elements of $A$ algebraic over a countable domain
form a countable set (\texttt{Algebraic.countable}). Its preimage in $\mathbb R$ is countable,
and the complement of each point is open and dense, so the complement of the preimage is
residual (countable intersection).
\end{proof}

\begin{lemma}[\texttt{indep\_pair}]
If $x\in\mathbb R$ is transcendental over $\mathbb Q$ and $y$ is transcendental over
$\mathbb Q[x]$ (in $A$), then $x\ne y$ and $\{x,y\}\subseteq A$ is algebraically independent over
$\mathbb Q$.
\end{lemma}
\begin{proof}
$\{x\}$ is independent because $x$ is transcendental (\texttt{algebraicIndependent\_unique\_type\_iff}).
Adding an element transcendental over the algebra generated by the set keeps independence
(\texttt{AlgebraicIndepOn.insert}). And $y\ne x$ since $x\in\mathbb Q[x]$ is algebraic over it.
\end{proof}

\begin{theorem}[\texttt{exists\_two\_indep\_U}]
Let $x=\sum_{i\ge0}10^{-i!}$ (Mathlib \texttt{liouvilleNumber 10}; the sum starts at $i=0$, so
$x$ is $1/10$ plus Liouville's constant $\sum_{k\ge1}10^{-k!}$; Liouville by
\texttt{liouville\_liouvilleNumber}). There is a real $y\ne x$ such that $x,y$ are U-numbers and
$\{x,y\}$ is algebraically independent over $\mathbb Q$, both in $\mathbb R$ and in $\mathbb C$.
\end{theorem}
\begin{proof}
The Liouville numbers form a residual set (Mathlib \texttt{eventually\_residual\_liouville}).
Intersect it with the two residual sets of the previous lemma ($A=\mathbb R$ and
$A=\mathbb C$). The intersection is residual, hence dense by the Baire category theorem
(\texttt{dense\_of\_mem\_residual}) and non-empty; pick $y$ in it. Then $y$ is Liouville, hence a
U-number, and the previous lemma gives independence.
\end{proof}

\begin{theorem}[\texttt{not\_conjectureReal}, \texttt{not\_conjectureComplex},
\texttt{two\_le\_iSup\_real}, \texttt{two\_le\_iSup\_complex}]
The set $s=\{x,y\}$ is a two-element algebraically independent set of U-numbers, in $\mathbb R$
and in $\mathbb C$ (\texttt{two\_dim\_U\_family\_real}, \texttt{two\_dim\_U\_family\_complex}).
So (C3b) fails, the supremum in (C3a) is $\ge2$, and \texttt{ConjectureReal} and
\texttt{ConjectureComplex} are false.
\end{theorem}

\section{Lean formalization}

File \texttt{lean/Conjecture9682/Basic.lean} (300 lines; only \texttt{import Mathlib}, Lean
v4.33.1 with Mathlib v4.33.1). Main theorems:
(Unicode symbols transliterated to ASCII.)
\begin{verbatim}
theorem not_conjectureComplex : Not ConjectureComplex
theorem not_conjectureReal : Not ConjectureReal
theorem two_le_iSup_complex : 2 <= iSup (fun s : {s : Set Complex //
    s subset {x | IsUNumber x} and
    AlgebraicIndependent Rat (Subtype.val : s -> Complex)} => s.1.encard)
theorem exists_two_indep_U : exists x y : Real, x != y and
    IsUNumber (x : Complex) and IsUNumber (y : Complex) and
    AlgebraicIndependent Rat (Subtype.val : ({x, y} : Set Real) -> Real) and
    AlgebraicIndependent Rat
      (Subtype.val : ({(x : Complex), (y : Complex)} : Set Complex) -> Complex)
\end{verbatim}
Axiom audit (\texttt{\#print axioms}) for \texttt{not\_conjectureComplex},
\texttt{not\_conjectureReal}, \texttt{two\_le\_iSup\_complex}, \texttt{two\_le\_iSup\_real} and
\texttt{liouville\_isUNumber}: \texttt{[propext, Classical.choice, Quot.sound]}. There is no
\texttt{sorry} and no \texttt{native\_decide}.

\end{document}
